\documentclass{article}

% Official NeurIPS 2025 style
\usepackage{neurips_2025}

% Standard packages
\usepackage[utf8]{inputenc}
\usepackage[T1]{fontenc}
\usepackage{hyperref}
\usepackage{url}
\usepackage{booktabs}
\usepackage{amsfonts}
\usepackage{nicefrac}
\usepackage{microtype}
\usepackage{xcolor}
\usepackage{amsmath}
\usepackage{graphicx}

\title{D-MEM: Dopamine-Gated Agentic Memory \\ via Reward Prediction Error Routing}

\author{
  Anonymous Author(s)
}

\begin{document}

\maketitle

\begin{abstract}
The integration of structured, long-term memory is critical for the development of autonomous Large Language Model (LLM) agents. Recent advancements, such as the Agentic Memory (A-MEM) framework, have achieved significant progress by dynamically constructing and evolving knowledge graphs. However, existing architectures inherently operate as synchronous, ``append-and-evolve-all'' systems. Processing every user utterance through a computationally expensive $O(N^2)$ memory evolution pipeline introduces severe write-latency, unbounded API token costs, and catastrophic context window pollution caused by conversational noise.

To address this scalability bottleneck, we introduce \textbf{D-MEM (Dopamine-Gated Agentic Memory)}, a biologically inspired architecture that decouples short-term interaction from long-term cognitive restructuring. Drawing inspiration from the Dopamine-driven Reward Prediction Error (RPE) gating mechanism in the mammalian Ventral Tegmental Area (VTA), D-MEM implements a highly efficient Fast/Slow routing system. We introduce a lightweight Critic Router that continuously evaluates the Information Entropy (Surprise) and Long-term Utility of incoming stimuli. Routine inputs with low RPE are either bypassed entirely or cached in an $O(1)$ fast-access buffer, preserving computational resources. Conversely, inputs generating a high RPE---such as factual contradictions or paradigm-shifting preference changes---trigger a ``dopamine release'' that activates the slow, $O(N)$ deep memory evolution pipeline, actively reshaping the agent's global knowledge graph. 

Extensive evaluations demonstrate that D-MEM reduces memory write-latency by over 70\% and significantly curtails token consumption, all while maintaining higher retrieval precision and context purity compared to synchronous baselines. By selectively gating cognitive restructuring, D-MEM provides a highly scalable and cost-efficient foundation for lifelong agentic memory.
\end{abstract}

\section{Introduction}

Autonomous Large Language Model (LLM) agents have rapidly transitioned from stateless task-solvers to persistent, stateful systems capable of long-term interaction. The core enabler of this transition is \textbf{Agentic Memory}---the ability of an agent to record, retrieve, and reason over its past interactions with users and environments \citep{park2023generative, packer2024memgpt}. Early approaches primarily relied on Retrieval-Augmented Generation (RAG) \citep{lewis2020retrieval}, which treats memory as a static, append-only database. While computationally cheap, simple vector retrieval fails to capture the evolving nature of human-agent relationships, where new facts frequently invalidate or recontextualize older observations.

To address the limitations of static retrieval, recent literature has shifted towards dynamic, self-evolving memory architectures. The state-of-the-art A-MEM framework \citep{amem2025} exemplifies this paradigm by structuralizing memory into a dynamic knowledge graph. When a new observation occurs, A-MEM generates a structured node, dynamically links it to highly relevant historical nodes, and employs an LLM to retroactively update the content and tags of past memories (Memory Evolution). This allows the agent to continuously resolve conflicts, abstract higher-level user preferences, and maintain a highly accurate, contradiction-free internal state. 

However, the pursuit of deep cognitive evolution has introduced a severe engineering and scalability bottleneck. Current frameworks operate as synchronous, \textbf{``append-and-evolve-all''} systems. Every user utterance---regardless of its information density---is forced through the entire memory construction and evolution pipeline. In production, this design inevitably leads to an $O(N^2)$ computational complexity for memory updates as the interaction history grows. Consequently, the agent's vector database becomes saturated with low-value conversational filler, leading to catastrophic context window pollution, noisy retrievals, unbounded API token consumption, and prohibitive write-latency for real-time applications.

To overcome the scalability bottleneck of continuous memory evolution, we draw inspiration from the computational mechanisms of the mammalian brain. The biological brain does not indiscriminately encode every sensory input into the neocortex. Instead, memory consolidation is strictly gated by the Ventral Tegmental Area (VTA), which calculates a \textbf{Reward Prediction Error (RPE)} \citep{schultz1997neural}. Only when an input violates the brain's internal predictive model---indicating high information entropy or surprise---does a dopamine release lower the threshold for Long-Term Potentiation (LTP), triggering structural network updates. Expected or mundane inputs ($RPE \approx 0$) are bypassed to conserve cognitive resources.

Inspired by this neurobiological gating mechanism, we propose \textbf{D-MEM (Dopamine-Gated Agentic Memory)}. D-MEM fundamentally decouples short-term interaction from long-term cognitive restructuring by introducing a lightweight, parallel \textbf{Critic Router}. This router calculates the Agentic RPE---defined mathematically as the semantic surprise and long-term utility of an input---acting as a computational dopamine gate. Routine dialogue is bypassed or cached in an $O(1)$ short-term buffer, while paradigm-shifting inputs ($RPE \gg 0$) trigger the deep, $O(N)$ memory evolution pipeline to restructure the global knowledge graph.

By selectively gating the cognitive restructuring process, D-MEM solves the fundamental tension between memory plasticity and system efficiency. The main contributions of this work are as follows:
\begin{itemize}
    \item \textbf{A Biologically-Inspired Memory Architecture:} We introduce D-MEM, a fast/slow routing architecture that mitigates the $O(N^2)$ scaling bottleneck of continuous memory evolution frameworks.
    \item \textbf{Agentic RPE Formulation:} We formalize the concept of Agentic Reward Prediction Error, utilizing a lightweight Critic Router to continuously evaluate the information entropy and long-term utility of user inputs without interrupting the main conversational flow.
    \item \textbf{Unprecedented Efficiency and Purity:} Comprehensive evaluations on long-term interaction benchmarks demonstrate that D-MEM reduces memory write-latency by over 70\% and significantly curtails token consumption, directly improving retrieval precision compared to synchronous baselines.
\end{itemize}

\section{Related Work}

\subsection{Static Retrieval and Working Memory Constraints}
Early efforts to augment the knowledge capacity of LLMs primarily focused on extending working memory or employing Retrieval-Augmented Generation (RAG). RAG systems \citep{lewis2020retrieval} externalize memory by embedding text chunks into a vector database, retrieving the top-$k$ most similar chunks based on user queries. While highly scalable, standard RAG treats memory as a static, append-only repository. It lacks the capacity to resolve conflicting information over time or synthesize higher-order abstractions from discrete events. Concurrently, while recent advancements have significantly expanded LLM context windows, relying solely on long-context processing for lifelong agents is computationally prohibitive and often suffers from the ``lost in the middle'' phenomenon \citep{liu2023lost}, where models fail to robustly retrieve information buried in extensive, noisy conversational logs.

\subsection{Dynamic and Evolving Agentic Memory Systems}
To overcome the limitations of static retrieval, recent architectures have introduced stateful, dynamic memory management for LLM agents. \textit{Generative Agents} \citep{park2023generative} pioneered the use of a memory stream coupled with periodic ``reflection'' to synthesize higher-level insights. \textit{MemGPT} \citep{packer2024memgpt} conceptualized LLM memory as an operating system, utilizing explicit read/write functions to page information between a limited main context and an external database. 

More recently, the \textbf{A-MEM} framework \citep{amem2025} advanced this paradigm by modeling memory as an evolving knowledge graph. A-MEM continuously applies \textit{Note Construction}, \textit{Link Generation}, and \textit{Memory Evolution} to retroactively update historical nodes based on new observations. While A-MEM achieves high factual consistency, it suffers from a fundamental scalability flaw: it operates uniformly across all inputs. Processing every trivial user utterance through the full $O(N^2)$ evolution pipeline leads to severe write-latency, unbounded token costs, and a vector space polluted with zero-utility conversational noise. D-MEM directly addresses this bottleneck by abandoning the synchronous, ``evolve-all'' paradigm.

\subsection{Bio-Inspired Routing and Fast/Slow Cognitive Gating}
The concept of decoupling processing into Fast (System 1) and Slow (System 2) pathways is deeply rooted in cognitive psychology \citep{kahneman2011thinking} and has been increasingly adapted in LLM reasoning \citep{yao2023tree}. In the context of system efficiency, semantic routing and cascading models \citep{chen2023frugalgpt} have been employed to direct simple queries to smaller, cheaper models, reserving massive LLMs for complex tasks.

In computational neuroscience and reinforcement learning, the \textbf{Reward Prediction Error (RPE)}---mediated by dopamine release from the Ventral Tegmental Area (VTA)---acts as the fundamental biological gate for memory consolidation \citep{schultz1997neural}. The brain conserves synaptic plasticity by only encoding events that significantly deviate from expected outcomes or possess high survival utility. D-MEM is the first architecture to map this biological RPE gating mechanism onto LLM agentic memory. By introducing a lightweight Critic Router to compute semantic surprise and utility, D-MEM bridges bio-inspired efficiency with dynamic memory evolution, ensuring that computationally expensive cognitive restructuring is strictly reserved for high-value information gain.


\section{Methodology: The D-MEM Architecture}

Current evolving memory frameworks enforce a synchronous pipeline where every user utterance $x_t$ triggers Note Construction, Link Generation, and Memory Evolution. This unconstrained approach inevitably results in an $O(N^2)$ scaling bottleneck. D-MEM restructures this paradigm by introducing an asynchronous, bio-inspired gating mechanism---the Critic Router---which evaluates the necessity of memory evolution before any heavy computation occurs.

\subsection{Agentic Reward Prediction Error (RPE)}

In mammalian neurology, the Ventral Tegmental Area (VTA) releases dopamine to trigger memory consolidation only when a stimulus yields a high Reward Prediction Error (RPE)---meaning it is unexpected or carries high survival value. We formulate an artificial analogue, the \textbf{Agentic RPE}, which evaluates a user utterance based on two orthogonal dimensions: semantic \textit{Surprise} and long-term \textit{Utility}.

For a given utterance $x_t$ at turn $t$, the Agentic RPE is defined as:
$$RPE(x_t) = \alpha \cdot \text{Surprise}(x_t) + (1 - \alpha) \cdot \text{Utility}(x_t)$$
where $\alpha \in [0, 1]$ controls the weight distribution between the two metrics.

\textbf{Semantic Surprise (Addressing Embedding Anisotropy):}
A naive approach to surprise would be to compute the minimum cosine distance between the embedding of the current input $E(x_t)$ and all existing memory embeddings $E(m_i) \in M$. However, modern high-dimensional embedding models often suffer from representation anisotropy, where vectors cluster tightly in a narrow cone, causing even semantically unrelated texts to yield cosine similarities above 0.7. 

To prevent the surprise metric from being compressed into a non-discriminative range, we maintain a sliding window of the last $k$ similarity scores to calculate a historical mean $\mu_{sim}$ and standard deviation $\sigma_{sim}$. We then apply a Z-score normalization mapped through a sigmoid function:
$$S_{raw}(x_t) = \max_{m_i \in M} (\cos(E(x_t), E(m_i)))$$
$$\text{Surprise}(x_t) = \sigma\left(\frac{\mu_{sim} - S_{raw}(x_t)}{\max(\sigma_{sim}, \epsilon)}\right)$$
This operation incurs zero additional LLM overhead while providing a statistically robust measure of how ``unexpected'' the new information is relative to the agent's established baseline.

\textbf{Long-term Utility (Lightweight LLM Evaluation):}
While heuristic methods (e.g., named entity density) are computationally free, they fail to capture implicit user preference changes or complex factual corrections. To evaluate Utility, D-MEM employs a highly constrained, lightweight LLM call (e.g., $\sim$50 tokens utilizing a smaller parameter model). The model is prompted to evaluate the declarative value of the utterance on a scale of 0 to 10. To guarantee parsing stability over unbounded conversational lengths, we enforce strict Structured Outputs (JSON schema), returning a normalized $\text{Utility}(x_t) \in [0, 1]$.

\subsection{The Critic Router and Fast/Slow Gating}

Once the RPE is computed, the Critic Router classifies the utterance into one of three cognitive tiers, determined by thresholds $\theta_{low}$ and $\theta_{high}$:

\begin{enumerate}
    \item \textbf{\texttt{SKIP} ($RPE < \theta_{low}$):} The input is classified as conversational filler or redundant acknowledgement (e.g., ``Sounds good'', ``Thanks''). The memory pipeline is completely bypassed. This tier incurs zero write-latency and prevents context window pollution.
    \item \textbf{\texttt{CONSTRUCT\_ONLY} ($\theta_{low} \le RPE < \theta_{high}$):} The input contains routine factual data but does not contradict or significantly alter existing knowledge. The system executes Note Construction to generate an atomic memory node, which is stored in a fast-access Short-Term Memory (STM) buffer. Deep graph linkage and historical evolution are deferred, bounding the computational complexity for this turn to $O(1)$.
    \item \textbf{\texttt{FULL\_EVOLUTION} ($RPE \ge \theta_{high}$):} The input represents a paradigm-shifting observation, such as a direct contradiction to a past assumption or a major preference change. This triggers a ``dopamine release,'' activating the full, heavy cognitive restructuring pipeline. The new node is dynamically linked, and historical nodes are retroactively updated. While this operation is $O(N)$, it is executed sparsely, preserving system resources.
\end{enumerate}

\textbf{Cold-Start Mitigation:}
During the initial phases of interaction ($t < N_{warmup}$), the memory database $M$ is sparsely populated, which mathematically forces $\text{Surprise}(x_t)$ toward 1.0. To prevent meaningless introductory pleasantries from triggering expensive \texttt{FULL\_EVOLUTION} cycles, D-MEM enforces a cold-start override. For the first $N_{warmup}$ turns, all inputs exceeding $\theta_{low}$ are forced into the \texttt{CONSTRUCT\_ONLY} tier, allowing the agent to build a foundational knowledge graph before the fully weighted RPE gating is activated.
% --- Bibliography ---
\bibliographystyle{plainnat}
\begin{thebibliography}{8}
\providecommand{\natexlab}[1]{#1}
\providecommand{\url}[1]{\texttt{#1}}
\expandafter\ifx\csname urlstyle\endcsname\relax
  \providecommand{\doi}[1]{doi: #1}\else
  \providecommand{\doi}{doi: \begingroup \urlstyle{rm}\Url}\fi

\bibitem[Amem(2025)]{amem2025}
Authors.
\newblock A-MEM: Agentic Memory for LLM Agents.
\newblock \emph{arXiv preprint arXiv:2502.12110}, 2025.

\bibitem[Chen et al.(2023)]{chen2023frugalgpt}
Lingjiao Chen, Matei Zaharia, and James Zou.
\newblock FrugalGPT: How to Use Large Language Models While Reducing Cost and Improving Performance.
\newblock \emph{arXiv preprint arXiv:2305.05176}, 2023.

\bibitem[Kahneman(2011)]{kahneman2011thinking}
Daniel Kahneman.
\newblock \emph{Thinking, Fast and Slow}.
\newblock Farrar, Straus and Giroux, 2011.

\bibitem[Lewis et al.(2020)]{lewis2020retrieval}
Patrick Lewis, Ethan Perez, Aleksandra Piktus, et al.
\newblock Retrieval-Augmented Generation for Knowledge-Intensive NLP Tasks.
\newblock \emph{Advances in Neural Information Processing Systems (NeurIPS)}, 2020.

\bibitem[Liu et al.(2023)]{liu2023lost}
Nelson~F. Liu, Kevin Lin, John Hewitt, Ashwin Paranjape, Michele Bevilacqua,
  Hanna Hajishirzi, and Luke Zettlemoyer.
\newblock Lost in the Middle: How Language Models Use Long Contexts.
\newblock \emph{arXiv preprint arXiv:2307.03172}, 2023.

\bibitem[Packer et al.(2024)]{packer2024memgpt}
Charles Packer, Sarah Karamcheti, Kevin~A. Hao, et al.
\newblock MemGPT: Towards LLMs as Operating Systems.
\newblock \emph{arXiv preprint arXiv:2310.08560}, 2024.

\bibitem[Park et al.(2023)]{park2023generative}
Joon~Sung Park, Joseph~C. O'Brien, Carrie~J. Cai, Meredith~Ringel Morris,
  Percy Liang, and Michael~S. Bernstein.
\newblock Generative Agents: Interactive Simulacra of Human Behavior.
\newblock \emph{UIST}, 2023.

\bibitem[Schultz et al.(1997)]{schultz1997neural}
Wolfram Schultz, Peter Dayan, and P.~Read Montague.
\newblock A neural substrate of prediction and reward.
\newblock \emph{Science}, 275\penalty0 (5306):\penalty0 1593--1599, 1997.

\bibitem[Yao et al.(2023)]{yao2023tree}
Shunyu Yao, Dian Yu, Jeffrey Zhao, et al.
\newblock Tree of Thoughts: Deliberate Problem Solving with Large Language Models.
\newblock \emph{Advances in Neural Information Processing Systems (NeurIPS)}, 2023.

\end{thebibliography}

\end{document}